% 解説プリントのコメント（saiten が書く。./tools/kaisetsu.ps1 がこのファイルからプリントを作る）
% 2026-11-22 のスキャン（p2・p3・p8・p9）の採点から作った
%% matome
\begin{daiji}
  \item 井戸の中の節は，(3) で決まる $\xi$ の範囲 $(n-1)\pi/2\le\xi<n\pi/2$ の両端を使って数える。上の端だけでは「多くても $n-1$ 個」しか言えない。奇パリティは $x=0$ の節が加わるので「同様」で済ませない（問題3 (4)）。
  \item 「$n$ 番目の状態が存在する条件」$(n-1)\pi/2<\xi_{0}$ と，「束縛状態がちょうど $n$ 個ある条件」$(n-1)\pi/2<\xi_{0}<n\pi/2$ を区別する（問題3 (3)）。
\end{daiji}
\begin{maemo}
  \item 数え上げの根拠の抜け：第11章 問題4・問題7（2026-11-20）で $n$ を正の整数に限る理由が抜けた。番号や個数を数えるときは，範囲の両端（どこから・どこまで）の根拠を書く。
\end{maemo}
%% mondai 3
\begin{kotaeran}
  \machigai{3}{○}{記述} 円の半径を大きくすると枝との交点が偶・奇の順に1つずつ増えること，基底状態が偶であること，$N=[2\xi_{0}/\pi+1]$（$2\xi_{0}/\pi$ が整数でなければ $\lceil 2\xi_{0}/\pi\rceil$ と同じ）は正しい。ただし「$n$ 番目の状態が存在する条件」として $n\pi/2>\xi_{0}\ge(n-1)\pi/2$ と書いた。これは束縛状態がちょうど $n$ 個ある条件で，$n$ 番目が存在する条件は $(n-1)\pi/2<\xi_{0}$ だけ（上限は要らない）。また，交点の順番がそのままエネルギーの順になる根拠（各区間に状態は高々1つで，$E=\hbar^{2}\xi^{2}/(2ma^{2})$ は $\xi$ について増える）が無い。
  \machigai{4}{△}{記述} 基底状態・第1励起状態の概形と，井戸の外は指数関数なので節が無いことは正しい。偶の状態の節を $\abs{(l-1/2)\pi}<\xi<\xi_{0}<n\pi/2$ から数えて $n-1$ 個としたが，(a) 使ったのは上の端 $\xi<n\pi/2$ だけで，これでは節は多くても $n-1$ 個としか言えない（$\xi\ge(n-1)\pi/2$ を使って，候補の $n-1$ 個がすべて $\abs{x}<a$ に入ることを言う必要がある），(b) $n$ 番目の状態について $\xi_{0}<n\pi/2$ とは限らない（それより上の状態もあり得る。言えるのは $\xi<n\pi/2$），(c) 奇パリティを「同様」としたが，奇では $x=0$ の節が加わり数え方が違う。
\end{kotaeran}
\begin{rikai}
  \item 偶（$n=2l+1$，$l\pi\le\xi<l\pi+\pi/2$）では，$\cos(\xi x/a)=0$ となる $\xi\abs{x}/a=\pi/2,\,3\pi/2,\,\dots$ のうち $\xi$ より小さいものが片側に $l$ 個で，両側で $2l=n-1$ 個。奇（$n=2l+2$，$l\pi+\pi/2\le\xi<(l+1)\pi$）では，$\sin(\xi x/a)=0$ の点が $x=0$ と片側 $l$ 個ずつで $2l+1=n-1$ 個。
  \item 「ちょうど何個」を示すには，上からの評価（多くても）と下からの評価（少なくとも）の両方が要る。
  \item 準位の順番：区間 $(n-1)\pi/2\le\xi<n\pi/2$ は左から偶・奇・偶…の枝が占め，$E\propto\xi^{2}$ なので区間の並びがそのままエネルギーの順になる。状態の有無は区間の左端と $\xi_{0}$ の大小だけで決まる。
\end{rikai}
